% SPDX-License-Identifier: Apache-2.0
\section{A Compound Value in a Waveform}
\label{sec:x-packet}

A packet is a struct of four fields, one of them an enum, and a
transaction that moves over a channel. In a trace a compound value is
one bit vector, its fields concatenated, and beside it one signal per
field under the value's name: the \code{Value} derive supplies the
parts, and a register or a channel that holds the value registers them
too. A viewer then shows \code{src}, \code{dst}, \code{kind} and
\code{len} as themselves, with the enum as its variant's index, and
Figure~\ref{fig:x-packet} shows exactly that. The file is FST when
\code{TXHDL\_FST} names one, which is what the document build asks
for now: \code{Wave::from\_env} picks FST or VCD by the variable set,
and \code{vcdcvt} reads either by the file's extension. The FST writer
used has bit vectors and reals and no string variables, so a value
keeps its bits rather than a name; that is the limit the per-field
signals work around.

\begin{figure}[htbp]
\centering
\input{packet_timing.tex}
\caption{The packet's fields, each a signal of its own, and the listener's copy of the kind.}
\label{fig:x-packet}
\end{figure}

\lstinputlisting[caption={\code{ex\_packet.rs}. Run with \code{bazel run //lib/examples:ex\_packet}.},label={lst:x-packet}]{lib/examples/ex_packet.rs}

\lstinputlisting[style=ascii,caption={What \code{ex\_packet} printed when the build ran it.},label={lst:x-packet-out}]{out_packet.txt}
